% ═══════════════════════════════════════════════════════════════════════
% The Geometry of Prime Vacuums: 
% Legendre's Conjecture in Function Fields via Monodromy and Chebotarev Density
% Titan Project — February 2026
% ═══════════════════════════════════════════════════════════════════════

\documentclass[11pt, a4paper]{article}

% ── Packages ──
\usepackage[top=28mm, bottom=28mm, left=25mm, right=25mm]{geometry}
\usepackage[T1]{fontenc}
\usepackage{amsmath, amssymb, amsthm, mathtools}
\usepackage{mathrsfs}
\usepackage{graphicx}
\usepackage[dvipsnames]{xcolor}
\usepackage{enumitem}
\usepackage{hyperref}

% ── Theorem Environments ──
\newtheorem{theorem}{Theorem}[section]
\newtheorem{lemma}[theorem]{Lemma}
\newtheorem{proposition}[theorem]{Proposition}
\newtheorem{corollary}[theorem]{Corollary}
\newtheorem{conjecture}[theorem]{Conjecture}
\theoremstyle{definition}
\newtheorem{definition}[theorem]{Definition}
\newtheorem{example}[theorem]{Example}
\theoremstyle{remark}
\newtheorem{remark}[theorem]{Remark}

% ── Macros ──
\newcommand{\Fq}{\mathbb{F}_q}
\newcommand{\Fqbar}{\overline{\mathbb{F}}_q}
\newcommand{\Proj}{\mathbb{P}^1}
\newcommand{\Affine}{\mathbb{A}}
\newcommand{\Gal}{\mathrm{Gal}}
\newcommand{\Frob}{\mathrm{Frob}}
\newcommand{\Sw}{\mathrm{Sw}}
\newcommand{\drop}{\mathrm{drop}}
\newcommand{\disc}{\mathrm{disc}}

\title{\textbf{The Geometry of Prime Vacuums: \\
Legendre's Conjecture in Function Fields \\
via Monodromy and Chebotarev Density}}
\author{Ruqing Chen\\[4pt]
\small GUT Geoservice Inc., Montr\'eal, Canada\\[2pt]
\small \texttt{ruqing@hotmail.com}}
\date{February 2026}

\begin{document}

\maketitle

\begin{abstract}
We prove the function field analogue of Legendre's Conjecture.
Let $f(t) \in \Fq[t]$ be a squarefree monic polynomial of degree
$d \ge 1$ with $\mathrm{char}(\Fq) > 2d$.
We show that for $q > (8d)^{2d+6}$,
the short interval
$\mathcal{I}_f = \{f(t)^2 + s(t) : \deg s \le d\}$
contains at least one irreducible polynomial over $\Fq$;
in fact, the number of such irreducibles is asymptotically
$q^{d+1}/(2d)$.
The proof constructs the universal polynomial family over the
parameter space $\Affine^{d+1}$, establishes that its geometric
monodromy group is the full symmetric group $S_{2d}$
(via irreducibility, primitivity, and simple branching), and
then applies the effective Chebotarev density theorem for
higher-dimensional varieties.
The topological error term is controlled by Katz's Betti
number bounds.
\end{abstract}

\bigskip

% ═══════════════════════════════════════════════════════════════════════
\section{Introduction}\label{sec:intro}
% ═══════════════════════════════════════════════════════════════════════

\subsection{Background}
The classical Legendre Conjecture posits that for every integer
$n \ge 1$, there exists at least one prime in the interval
$[n^2,(n+1)^2]$.
This remains open over $\mathbb{Z}$; the best unconditional result
(Baker--Harman--Pintz \cite{BHP2001}) yields primes in
$[x, x + x^{0.525}]$, which falls short of the
$\mathcal{O}(x^{0.5})$ length of the Legendre interval.

Over the polynomial ring $\Fq[t]$, the geometry of algebraic
varieties over finite fields provides the additional structure
needed to resolve the problem.
Our method follows the paradigm introduced by
Bary-Soroker \cite{barysoroker2012}, who proved the function field
analogue of Dirichlet's theorem on primes in arithmetic progressions
using the monodromy of polynomial families and the Chebotarev density
theorem.
We adapt and extend this framework to the substantially different
setting of \emph{short intervals}, where the shift $s(t)$ perturbs
only the lower-order coefficients of $f(t)^2$, introducing rigidity
constraints that require a more delicate analysis of primitivity
and branching.

\subsection{Context within the Titan Project}

This paper is part of a broader programme investigating the
geometric and cohomological obstructions to classical prime
number conjectures.
In \cite{ChenConductor}, we studied conductor rigidity phenomena
for Frey curves associated to prime gaps,
establishing that the Wiles paradigm encounters fundamental
obstructions in the additive setting;
this work provided the original motivation for seeking
alternative geometric approaches to Legendre-type problems.
In \cite{ChenLandau}, we investigated Landau's fourth problem
($n^2 + 1$ primes) via conductor rigidity and Sato--Tate
equidistribution for quadratic polynomial families---a setting
closely analogous to the present one, where the quadratic
structure $f(t)^2 + s(t)$ plays a similar role.
In \cite{ChenAP}, we examined primes in arithmetic progressions
through the lens of conductor rigidity,
encountering the Bombieri--Vinogradov barrier;
the monodromy and Chebotarev density framework employed in the
present paper may be viewed as a function field resolution of
the equidistribution difficulties identified there.

The present work differs from the preceding papers in a fundamental
way: rather than analyzing conductor invariants of fixed
$L$-functions, we construct a \emph{universal polynomial family}
and exploit the full symmetric monodromy of the associated covering
space.
This shift from local rigidity to global equidistribution is what
enables a complete resolution in the function field setting.

\subsection{Setup and main result}

\begin{definition}[Legendre Short Interval]\label{def:interval}
Let $f(t) \in \Fq[t]$ be monic of degree $d \ge 1$.
The \emph{Legendre short interval} is
\[
    \mathcal{I}_f
    \;=\;
    \Bigl\{\,
    P_\mathbf{a}(t) = f(t)^2 + \sum_{i=0}^d a_i\, t^i
    \;\Bigm|\; a_i \in \Fq
    \,\Bigr\}.
\]
The parameter space is $\mathcal{M} = \Affine^{d+1}_{\Fq}$
with coordinates $\mathbf{a}=(a_0,\dots,a_d)$.
Every $P_\mathbf{a}$ is monic of degree $2d$, and
$|\mathcal{I}_f| = q^{d+1}$.
\end{definition}

\begin{theorem}[Main Theorem]\label{thm:main}
Let $f(t) \in \Fq[t]$ be a squarefree monic polynomial of degree
$d \ge 1$, and suppose $\mathrm{char}(\Fq) > 2d$.
Then the number of irreducible polynomials in $\mathcal{I}_f$
satisfies
\begin{equation}\label{eq:main_asymptotic}
    N_{\mathrm{irr}}
    \;=\;
    \frac{q^{d+1}}{2d}
    \;+\;
    O_d\!\left(q^{d+\frac{1}{2}}\right).
\end{equation}
In particular, $N_{\mathrm{irr}} \ge 1$ whenever
$q > (8d)^{2d+6}$.
\end{theorem}

\begin{remark}[On the squarefree hypothesis]\label{rem:sqfree}
The assumption that $f$ is squarefree is used in
Proposition~\ref{prop:transposition}
to ensure that simple branching points
exist on the discriminant locus.
When $f$ has repeated roots, $f(t)^2$ acquires roots of multiplicity
$\ge 4$, and the analysis of local monodromy at such higher-order
singularities requires versal deformations of
$A_k$-singularities \cite{AGLV1993}.
We expect Theorem~\ref{thm:main} to hold without the squarefree
condition, but we impose it here for clarity.
\end{remark}

The proof occupies Sections~\ref{sec:monodromy}--\ref{sec:proof}
and proceeds in three steps:
\begin{enumerate}[nosep, label=(\roman*)]
\item Construct the universal polynomial family and prove its
      geometric monodromy group is $S_{2d}$
      (Section~\ref{sec:monodromy}).
\item Apply the effective Chebotarev density theorem for
      higher-dimensional parameter spaces
      (Section~\ref{sec:chebotarev}).
\item Bound the topological error term via Katz's Betti number
      estimates, complemented by the Grothendieck--Ogg--Shafarevich
      formula on one-dimensional slices (Section~\ref{sec:GOS}).
\end{enumerate}

% ═══════════════════════════════════════════════════════════════════════
\section{The Universal Family and Its Monodromy Group}
\label{sec:monodromy}
% ═══════════════════════════════════════════════════════════════════════

\subsection{The universal polynomial}\label{subsec:universal}

Fix a squarefree monic polynomial
$f(t) = t^d + c_{d-1}t^{d-1} + \cdots + c_0 \in \Fq[t]$.
Define the \emph{universal polynomial}
\begin{equation}\label{eq:universal}
    F(t,\mathbf{a})
    \;=\;
    f(t)^2 + a_0 + a_1\,t + \cdots + a_d\,t^d
    \;\in\; \Fq[\mathbf{a}][t],
\end{equation}
a monic polynomial of degree $2d$ in $t$.
Let $\mathcal{H} = \{F(t,\mathbf{a})=0\}
\subset \mathcal{M} \times \Affine^1_t$,
and $\pi: \mathcal{H} \to \mathcal{M}$ the projection.
This is a finite morphism of degree $2d$.

The coefficient of $t^k$ in $F$ is:
\begin{equation}\label{eq:coefficients}
    [t^k]\, F \;=\;
    \begin{cases}
    [t^k]\, f(t)^2 & \text{if } d+1 \le k \le 2d, \\[3pt]
    [t^k]\, f(t)^2 \;+\; a_k & \text{if } 0 \le k \le d.
    \end{cases}
\end{equation}
The top $d-1$ sub-leading coefficients
(of $t^{2d-1}, \ldots, t^{d+1}$) are fixed constants
determined by $f$, while the bottom $d+1$ coefficients
(of $t^d, \ldots, t^0$) vary freely with $\mathbf{a}$.
This \emph{split structure} is crucial for the primitivity argument.

\begin{definition}\label{def:disc_locus}
The \emph{discriminant locus} is
$\mathcal{D}
= \{\mathbf{a} \in \mathcal{M} :
\disc_t(F(t,\mathbf{a})) = 0\}$.
Over $\mathcal{M}^\circ = \mathcal{M} \setminus \mathcal{D}$,
the morphism $\pi$ is \'etale, defining the
\emph{geometric monodromy group}
\[
    G_{\mathrm{geom}}
    \;=\;
    \mathrm{Im}\!\left(\,
    \pi_1^{\mathrm{\acute{e}t}}
    \bigl(\mathcal{M}^\circ_{\Fqbar}\bigr)
    \;\to\; S_{2d}
    \,\right).
\]
\end{definition}

We establish $G_{\mathrm{geom}} = S_{2d}$ in three steps.

\subsection{Step 1: Irreducibility (Transitivity)}
\label{subsec:irred}

\begin{proposition}\label{prop:irred}
The polynomial $F(t,\mathbf{a})$ is irreducible over
$\Fqbar(\mathbf{a})$.
\end{proposition}

\begin{proof}
Write $F$ as an element of $R[a_0]$ where
$R = \Fqbar[t, a_1, \ldots, a_d]$:
\[
    F \;=\; a_0 + \bigl(f(t)^2 + a_1 t + \cdots + a_d t^d\bigr).
\]
Since $F$ is \emph{degree one} in $a_0$ with coefficients in the
integral domain $R$, it is irreducible in
$R[a_0] = \Fqbar[t, a_0, \ldots, a_d]$.

This ring is a UFD (polynomial ring over a field).
By Gauss's lemma, if $F$ were reducible in
$\Fqbar(\mathbf{a})[t]$, clearing denominators would yield
a nontrivial factorization in $\Fqbar[\mathbf{a}, t]$,
contradicting the above.
\end{proof}

\begin{corollary}\label{cor:transitive}
$G_{\mathrm{geom}}$ acts transitively on the $2d$ roots of $F$.
\end{corollary}

\begin{proof}
Irreducibility over $\Fqbar(\mathbf{a})$ implies the generic fiber
of $\mathcal{H}^\circ \to \mathcal{M}^\circ$ is connected,
equivalent to transitivity (\cite[Exp.~V, Prop.~8.2]{SGA1}).
\end{proof}

\subsection{Step 2: Primitivity}\label{subsec:primitive}

A transitive subgroup $G \le S_n$ is \emph{primitive} if it
preserves no nontrivial block partition of $\{1,\ldots,n\}$.
For polynomial coverings, imprimitivity is equivalent to a
functional decomposition $F(t) = g(h(t))$ with
$1 < \deg h < n$
(\cite{mueller1993, turnwald1995}).

\begin{proposition}\label{prop:primitive}
$G_{\mathrm{geom}}$ is primitive.
\end{proposition}

\begin{proof}
Suppose for contradiction that there exist monic polynomials
$h \in \Fqbar(\mathbf{a})[t]$ of degree $m$ and
$g \in \Fqbar(\mathbf{a})[x]$ of degree $n = 2d/m$
(with $1 < m < 2d$, $m \mid 2d$) such that
\begin{equation}\label{eq:decomp}
    F(t,\mathbf{a})
    \;=\; g\bigl(h(t,\mathbf{a}),\;\mathbf{a}\bigr).
\end{equation}
Since $m \mid 2d$ and $m < 2d$, we have $m \le d$ and $n \ge 2$.

\textbf{Normalization.}
Replacing $h(t)$ by $h(t) - h(0)$ and $g(x)$ by $g(x + h(0))$
preserves \eqref{eq:decomp} while achieving
\begin{equation}\label{eq:normalize}
    h(0) = 0.
\end{equation}
Write $h(t) = t^m + h_{m-1}t^{m-1} + \cdots + h_1\, t$
(with $h_i \in \Fqbar(\mathbf{a})$, no constant term), and
$g(x) = x^n + g_{n-1}x^{n-1} + \cdots + g_0$.

\textbf{Coefficient analysis.}
For $1 \le j \le m-1$, the coefficient $[t^{2d-j}]$ in the
expansion of $g(h(t))$ depends only on $h(t)^n$ at this degree
(since $g_{n-1} h(t)^{n-1}$ contributes at degree $\le 2d - m$,
and $2d - j > 2d - m$ for $j < m$).
Specifically:
\[
    [t^{2d-j}]\, g(h(t))
    \;=\; n\, h_{m-j}
    \;+\; \Phi_j(h_{m-1}, \ldots, h_{m-j+1})
\]
for a universal polynomial $\Phi_j$ (with $\Phi_1 = 0$).

Since $m \le d$, we have $2d - j \ge 2d - (m-1) \ge d + 1$
for all $1 \le j \le m - 1$.
By the split structure~\eqref{eq:coefficients},
$[t^{2d-j}]\, F$ is a \emph{fixed constant} in $\Fq$
for these degrees.

Equating inductively: for $j = 1$,
$n \cdot h_{m-1} = 2c_{d-1}$,
so $h_{m-1} = 2c_{d-1}/n \in \Fq$ (using
$\gcd(n, \mathrm{char}) = 1$ since $n \le 2d < p$).
For $j = 2, \ldots, m-1$, each step determines $h_{m-j}$
as a fixed element of $\Fq$ in terms of previously determined
coefficients.
Thus $h_{m-1}, \ldots, h_1 \in \Fq$ are all fixed constants.
Combined with the normalization $h_0 = 0$,
the polynomial $h(t) \in \Fq[t]$ is entirely independent of
$\mathbf{a}$.

\textbf{Contradiction.}
Since $h \in \Fq[t]$ is fixed, $F(t,\mathbf{a}) = g(h(t),\mathbf{a})$.
The coefficients of $F$ at degrees $0 \le k \le d$ are
$[t^k] f(t)^2 + a_k$, providing $d + 1$ independent free
parameters $(a_0, \ldots, a_d)$.
But $g$ has only $n = 2d/m$ free coefficients
$(g_0, \ldots, g_{n-1})$, and $n \le d < d + 1$.
The polynomial map
$(g_0, \ldots, g_{n-1})
\mapsto \bigl([t^0]\, g(h(t)), \ldots, [t^d]\, g(h(t))\bigr)$
goes from $\Affine^n$ to $\Affine^{d+1}$ with $n < d+1$,
hence is not surjective.
This contradicts the fact that the left-hand side ranges over
all of $\Affine^{d+1}$ as $\mathbf{a}$ varies.
\end{proof}

\begin{remark}\label{rem:d1_primitive}
When $d = 1$, the group $S_2$ has no proper transitive subgroup,
so primitivity is automatic.
\end{remark}

\subsection{Step 3: Transpositions from simple branching}
\label{subsec:transposition}

\begin{proposition}\label{prop:transposition}
Assume $f(t)$ is squarefree and $\mathrm{char}(\Fq) > 2d$.
Then $G_{\mathrm{geom}}$ contains a transposition.
\end{proposition}

\begin{proof}
We exhibit smooth points of $\mathcal{D}$ where exactly two roots
of $F$ coalesce.

\textbf{Setup.}
Since $f$ is squarefree, it has $d$ distinct roots
$\alpha_1, \ldots, \alpha_d \in \Fqbar$.
At $\mathbf{a} = \mathbf{0}$, the polynomial
$F(t, \mathbf{0}) = f(t)^2$ has each $\alpha_j$ as a double root.

\textbf{Independent unfolding via Vandermonde.}
Consider the evaluation map
$\varphi : \Fqbar^d \to \Fqbar^d$ defined by
\[
    \varphi(a_0, \ldots, a_{d-1})
    \;=\;
    \bigl(P(\alpha_1),\; P(\alpha_2),\;\ldots,\; P(\alpha_d)\bigr),
\]
where $P(t) = a_0 + a_1 t + \cdots + a_{d-1} t^{d-1}$.
The Jacobian of $\varphi$ is the Vandermonde matrix
$V(\alpha_1, \ldots, \alpha_d)$, whose determinant
$\prod_{i<j}(\alpha_j - \alpha_i)$
is nonzero since the $\alpha_j$ are distinct.
Therefore $\varphi$ is an isomorphism,
and the $d$ parameters $(a_0, \ldots, a_{d-1})$
independently control the perturbation at each double root.

Concretely, near $\alpha_j$ and $\mathbf{a} = \mathbf{0}$,
the local equation of $F$ in the coordinate $u = t - \alpha_j$ is
\[
    F(\alpha_j + u,\, \mathbf{a})
    \;=\;
    f'(\alpha_j)^2 u^2 + P(\alpha_j) + O(u \cdot \mathbf{a})
    + O(u^3).
\]
Since $f$ is squarefree, $f'(\alpha_j) \ne 0$
(using $\mathrm{char}(\Fq) > 2d \ge \deg f'+ 1$).
The double root at $u = 0$ splits into two simple roots when
$P(\alpha_j) \ne 0$, or coalesces when $P(\alpha_j) = 0$.
The invertibility of $\varphi$ means these $d$ conditions
are independent.

\textbf{Identifying simple nodes.}
The discriminant locus $\mathcal{D}$ near $\mathbf{a} = \mathbf{0}$
has $d$ smooth branches
$\mathcal{D}_j = \{\mathbf{a} : P(\alpha_j) = 0\}$
for $j = 1, \ldots, d$,
intersecting transversally at the origin.
A general point on $\mathcal{D}_j$
(away from the other branches) corresponds to a polynomial
$P_\mathbf{a}$ with exactly one double root (near $\alpha_j$)
and $2d - 2$ simple roots.

\textbf{Local monodromy.}
By the Picard--Lefschetz formula
(\cite[Exp.~XV, \S3]{SGA7}; \cite[\S I.6]{katz1988}),
an \'etale loop around the smooth locus of $\mathcal{D}_j$
transposes the two coalescing sheets and fixes all others.
This is a transposition in $G_{\mathrm{geom}}$.
\end{proof}

\begin{remark}\label{rem:nonsqfree}
When $f$ is not squarefree, double roots of $f$ become
fourfold roots of $f^2$, and local monodromy may produce
products of transpositions rather than a single transposition.
Extracting an individual transposition then requires
$A_k$-singularity analysis \cite{AGLV1993}.
We leave this extension to future work.
\end{remark}

\subsection{Conclusion: the full symmetric group}

\begin{theorem}\label{thm:monodromy}
Let $f \in \Fq[t]$ be squarefree monic of degree $d \ge 1$ with
$\mathrm{char}(\Fq) > 2d$.
Then $G_{\mathrm{geom}} = S_{2d}$.
\end{theorem}

\begin{proof}
$G_{\mathrm{geom}}$ is transitive
(Corollary~\ref{cor:transitive}),
primitive (Proposition~\ref{prop:primitive}),
and contains a transposition
(Proposition~\ref{prop:transposition}).
By Jordan's theorem
(\cite[Thm.~3.3E]{dixon1996}),
a \emph{primitive} subgroup of $S_n$ containing a transposition
equals $S_n$.
Hence $G_{\mathrm{geom}} = S_{2d}$.
\end{proof}

\begin{remark}\label{rem:jordan}
The hypothesis of \emph{primitivity} is essential.
The weaker statement
``transitive + contains a transposition $\Rightarrow S_n$''
is false: the dihedral group $D_4 \le S_4$ is transitive
and contains transpositions, but $D_4 \ne S_4$.
\end{remark}

% ═══════════════════════════════════════════════════════════════════════
\section{Counting Irreducibles via the Chebotarev Density Theorem}
\label{sec:chebotarev}
% ═══════════════════════════════════════════════════════════════════════

\subsection{Frobenius elements and factorization type}

For $\mathbf{a} \in \mathcal{M}^\circ(\Fq)$,
the Frobenius element $\Frob_\mathbf{a}$ is a well-defined
conjugacy class in $S_{2d}$ whose cycle type
$(d_1, \ldots, d_k)$ encodes the degrees of the irreducible
factors of $P_\mathbf{a}$ over $\Fq$.
Thus $P_\mathbf{a}$ is irreducible $\Leftrightarrow$
$\Frob_\mathbf{a}$ is a $2d$-cycle.

\begin{lemma}\label{lem:cycle_fraction}
The proportion of $2d$-cycles in $S_{2d}$ is $1/(2d)$.
\end{lemma}

\begin{proof}
There are $(2d-1)!$ such cycles, so the proportion is
$(2d-1)!/(2d)! = 1/(2d)$.
\end{proof}

\subsection{The effective Chebotarev theorem
for higher-dimensional varieties}

We use the following result, which combines the
Grothendieck--Lefschetz trace formula with Deligne's purity
theorem \cite{deligne1980}.
The formulation via the \emph{permutation sheaf} (rather than
the Galois closure) avoids factorial-sized constants.

\begin{theorem}[Effective Chebotarev;
cf.\ {\cite[Thm.~9.2.6]{katz1988}},
{\cite[\S9.7]{FJ2008}}]
\label{thm:chebotarev}
Let $U/\Fq$ be a smooth geometrically connected affine variety of
dimension $n$, and let $\pi : \mathcal{Y} \to U$ be a finite
\'etale cover of degree $N$ with connected geometric generic fiber,
so that the monodromy group $G \le S_N$ is transitive.
Let $\mathcal{F} = \pi_*\overline{\mathbb{Q}}_\ell$ be the
associated permutation sheaf of rank $N$ on $U$, and let
$C \subset G$ be a conjugacy class.
Then
\begin{equation}\label{eq:cheb}
    \left|\;
    \#\bigl\{\, u \in U(\Fq) : \Frob_u \in C \,\bigr\}
    \;-\; \frac{|C|}{|G|}\, q^n
    \;\right|
    \;\le\;
    B(\mathcal{F}) \cdot q^{n - 1/2}
    \;+\; O_{N,n}(q^{n-1}),
\end{equation}
where $B(\mathcal{F}) = \sum_{i=0}^{2n}
\dim H^i_c(U_{\Fqbar},\, \mathcal{F})$
is the total Betti number of the \emph{sheaf} $\mathcal{F}$
on $U$ (not the Betti number of the covering space).
\end{theorem}

\begin{remark}\label{rem:betti_bound}
The decisive advantage of formulating the error in terms of
$B(\mathcal{F})$ rather than the Betti numbers of the
Galois closure is that the former avoids factors of $|G|$
(which is $(2d)!$ for $G = S_{2d}$).
By \cite[Thm.~12]{katz2001}, for the complement $U$ of a
hypersurface of degree $\delta$ in $\Affine^{n}$:
\begin{equation}\label{eq:katz_bound}
    B(\mathcal{F}) \;\le\; N \cdot (2\delta + 1)^{n}.
\end{equation}
In our setting, $N = 2d$, $n = d + 1$, and $\delta$ is the
\emph{total degree} of $\disc_t(F)$ as a polynomial in
$\mathbf{a}$.
Since $F(t,\mathbf{a})$ is monic of degree $2d$ in $t$ and each
coefficient of $F$ is at most degree~$1$ in $\mathbf{a}$, the
discriminant --- being the resultant of $F$ and $F'$, hence
a determinant of size $(4d-1)$ whose entries are
coefficients of $F$ and $F'$ --- has total degree at most
\begin{equation}\label{eq:delta}
    \delta \;=\; 2(2d) - 2 \;=\; 4d - 2
\end{equation}
in $\mathbf{a}$ (this is the classical bound for the total degree of
the discriminant of a monic polynomial of degree $N$ in its
coefficients; it should \emph{not} be confused with the
\emph{weighted} degree $N(N-1) = 2d(2d-1)$, which assigns
weight $N - i$ to the $i$-th coefficient).

Hence $2\delta + 1 = 8d - 3$ and:
\begin{equation}\label{eq:B_explicit}
    B(\mathcal{F})
    \;\le\;
    2d \cdot (8d - 3)^{d+1}
    \;\le\;
    2d \cdot (8d)^{d+1}.
\end{equation}
The essential point: $B(\mathcal{F})$ depends on $d$ alone,
not on $q$.
\end{remark}

% ═══════════════════════════════════════════════════════════════════════
\section{Topological Complexity and the GOS Formula}
\label{sec:GOS}
% ═══════════════════════════════════════════════════════════════════════

While the Katz bound~\eqref{eq:B_explicit} suffices for the
asymptotic result, the Grothendieck--Ogg--Shafarevich (GOS)
formula provides sharper bounds on one-dimensional slices and
illuminates the geometric mechanism.

\subsection{The GOS formula on a generic pencil}

Let $\ell \cong \Affine^1 \hookrightarrow \mathcal{M}$ be a
general affine line.
After restricting and compactifying, we obtain a lisse
$\overline{\mathbb{Q}}_\ell$-sheaf $\mathcal{F}_\ell$
of rank $r = 2d - 1$
(the reduced permutation representation) on an open
$U_\ell \subset \Proj_{\Fqbar}$.
The GOS formula \cite{grothendieck1977, raynaud1965, serre1979}
states:
\begin{equation}\label{eq:GOS}
    \chi(\Proj_{\Fqbar},\, j_! \mathcal{F}_\ell)
    \;=\; 2r
    \;-\; \sum_{v \in \Proj \setminus U_\ell}
    \bigl(\drop_v + \Sw_v\bigr).
\end{equation}

\subsection{Tameness kills wild ramification}

Under $\mathrm{char}(\Fq) = p > 2d$:

\emph{At finite bad points $v$}: the inertia group
$I_v$ acts through a cyclic group of order dividing $2d$,
which is coprime to $p$.
Hence $\Sw_v(\mathcal{F}_\ell) = 0$
(\cite[IV, \S2, Cor.~1]{serre1979}).

\emph{At $v = \infty$}: the ramification indices divide $|S_{2d}|
= (2d)!$.
Since $p > 2d$, every integer from $1$ to $2d$ is coprime to $p$,
so the wild inertia subgroup (a $p$-group) acts trivially on the
$2d$ roots.
Thus $\Sw_\infty = 0$.

\subsection{One-dimensional Euler characteristic bound}

Since all Swan conductors vanish, \eqref{eq:GOS} gives
$\chi = 2r - \sum_v \drop_v$ with each
$0 \le \drop_v \le r = 2d - 1$.
By \eqref{eq:delta}, the number of finite bad fibers on a generic
pencil is at most $D = 4d - 2$, plus one point at $\infty$.
Hence:
\begin{equation}\label{eq:chi_1d}
    |\chi(\Proj_{\Fqbar}, j_!\mathcal{F}_\ell)|
    \;\le\; (D + 2) \cdot r
    \;\le\; 4d\,(2d-1)
    \;\le\; 8d^2.
\end{equation}

This provides a tight polynomial bound on one-dimensional slices.
The passage to the full $(d+1)$-dimensional parameter space
uses the Katz bound~\eqref{eq:B_explicit} as stated in
Remark~\ref{rem:betti_bound}.

% ═══════════════════════════════════════════════════════════════════════
\section{Proof of the Main Theorem}\label{sec:proof}
% ═══════════════════════════════════════════════════════════════════════

\begin{proof}[Proof of Theorem~\ref{thm:main}]
\textbf{Step 1: Monodromy.}
By Theorem~\ref{thm:monodromy}, $G_{\mathrm{geom}} = S_{2d}$.

\textbf{Step 2: Chebotarev count.}
Apply Theorem~\ref{thm:chebotarev} with $U = \mathcal{M}^\circ$,
$n = d + 1$, $N = 2d$, $C = $ conjugacy class of $2d$-cycles.
By Lemma~\ref{lem:cycle_fraction}, $|C|/|G| = 1/(2d)$.
Using the Betti bound~\eqref{eq:B_explicit}:
\begin{equation}\label{eq:count_intermediate}
    N_{\mathrm{irr}}^{\circ}
    \;:=\;
    \#\{\mathbf{a} \in \mathcal{M}^\circ(\Fq) :
    P_\mathbf{a} \text{ irreducible}\}
    \;=\;
    \frac{q^{d+1}}{2d}
    + O_d(q^{d+1/2}).
\end{equation}

\textbf{Step 3: Discriminant locus.}
$\mathcal{D}$ is a hypersurface of degree $\le 4d - 2$
in $\Affine^{d+1}$ (see \eqref{eq:delta}),
so $|\mathcal{D}(\Fq)| \le (4d-2)\,q^d + O_d(q^{d-1/2})$
by Lang--Weil \cite{LW1954}.
On $\mathcal{D}$, $P_\mathbf{a}$ has repeated roots, hence is
reducible (since irreducible implies separable when
$\mathrm{char} > \deg$).
Thus $N_{\mathrm{irr}} = N_{\mathrm{irr}}^\circ$.

\textbf{Step 4: Explicit positivity.}
Write
\begin{equation}\label{eq:positivity}
    N_{\mathrm{irr}}
    \;\ge\;
    \frac{q^{d+1}}{2d}
    - B(\mathcal{F}) \cdot q^{d+1/2}
    - C'_d \cdot q^d,
\end{equation}
where $B(\mathcal{F}) \le 2d\,(8d)^{d+1}$
(from \eqref{eq:B_explicit}) and
$C'_d \le 4d - 2$ absorbs the Lang--Weil term.
For $q \ge 2$, the $q^d$ term is dominated by $q^{d+1/2}$, so:
\[
    N_{\mathrm{irr}}
    \;\ge\;
    q^{d+1/2}\!\left(
    \frac{q^{1/2}}{2d} - B(\mathcal{F}) - C'_d
    \right).
\]
Since $C'_d \le 4d \le B(\mathcal{F})$ for $d \ge 1$,
we may take $A_d = 2\,B(\mathcal{F}) \le 4d\,(8d)^{d+1}$.
Then $N_{\mathrm{irr}} > 0$ when
$q^{1/2} > 2d \cdot A_d$, i.e.,
\[
    q \;>\; \bigl(2d \cdot A_d\bigr)^2
    \;=\; \bigl(8d^2\,(8d)^{d+1}\bigr)^2
    \;=\; 64\,d^4\,(8d)^{2d+2}.
\]
Since $64\,d^4 \le (8d)^4$ for $d \ge 1$, this yields:
\begin{equation}\label{eq:threshold}
    q \;>\; (8d)^{2d+6}.
\end{equation}
\end{proof}

\begin{remark}[On the threshold]\label{rem:threshold}
The bound \eqref{eq:threshold} is far from optimal;
it is driven by the worst-case Katz estimate for the Betti numbers.
The corrected discriminant degree $\delta = 4d - 2$ (as opposed to
the weighted degree $2d(2d-1)$ erroneously used in earlier versions)
yields a significant improvement: the base of the exponential is
$8d$ rather than $\sim 4d^2$.
Numerical evidence (Section~\ref{sec:numerical}) strongly suggests
that $N_{\mathrm{irr}} \ge 1$ holds for \emph{all} $q$ and $d$
with $\mathrm{char}(\Fq) > 2d$.
An approach via sieve methods or the large sieve inequality for
function fields (as in \cite{bank2015}) could potentially
yield much sharper bounds.
\end{remark}

% ═══════════════════════════════════════════════════════════════════════
\section{Numerical Verification}\label{sec:numerical}
% ═══════════════════════════════════════════════════════════════════════

\begin{example}[$d = 1$, $f(t) = t$]\label{ex:d1}
$\mathcal{I}_f$ is the set of all monic quadratics over $\Fq$.
The exact count of irreducibles is $(q^2-q)/2$,
matching $q^2/2 + O(q)$.

\smallskip
\noindent\textit{Data.}\;
$q = 5$: $N_{\mathrm{irr}} = 10$, predicted $12.5$.\quad
$q = 7$: $N_{\mathrm{irr}} = 21$, predicted $24.5$.\quad
$q = 11$: $N_{\mathrm{irr}} = 55$, predicted $60.5$.
\end{example}

\begin{example}[$d = 2$, $f(t) = t^2 + 1$, $q = 11$]
\label{ex:d2}
$f$ is squarefree (irreducible over $\mathbb{F}_{11}$).
$|\mathcal{I}_f| = 11^3 = 1331$;
predicted $\lfloor 1331/4 \rfloor = 332$;
direct computation gives $N_{\mathrm{irr}} = 330$.
\end{example}

\begin{example}[$d = 1$, $f(t) = t$, $q = 3$]
\label{ex:small}
$\mathrm{char} = 3 > 2 = 2d$, so the hypotheses hold.
The explicit asymptotic threshold is astronomically large,
yet $N_{\mathrm{irr}} = 3 > 0$.
This illustrates that the theorem's bound is sufficient
but vastly non-sharp.
\end{example}

% ═══════════════════════════════════════════════════════════════════════
\section{Discussion}\label{sec:discussion}
% ═══════════════════════════════════════════════════════════════════════

\subsection{The topological obstruction perspective}

The proof reduces to one principle:
\emph{a connected covering with full symmetric monodromy cannot
have all rational fibers reducible.}
$S_{2d}$ is the global symmetry; Chebotarev translates it into
equidistribution of Frobenius elements; the Katz/GOS bounds ensure
the error is $O(q^{-1/2})$ relative to the main term.

We emphasize that the GOS formula does \emph{not} produce
the contradiction directly (as was incorrectly attempted in
an earlier version via ``Swan conductor explosion'').
Swan conductors are local invariants of a single sheaf
and are always finite; the global factorization pattern is
an arithmetic property of the ensemble of rational fibers,
governed by equidistribution.

\subsection{Relation to prior work}

Bary-Soroker \cite{barysoroker2012} proved the function field
Dirichlet theorem by the same monodromy paradigm.
Our contribution is the adaptation to short intervals,
which differs in two respects:
(a) the parameter space is $(d+1)$-dimensional
(requiring the higher-dimensional Chebotarev theorem), and
(b) the rigid leading-term structure of $f(t)^2 + s(t)$
demands a non-trivial primitivity analysis
(Proposition~\ref{prop:primitive}).

\subsection{Relation to the classical Legendre Conjecture}

The function field proof succeeds because the Riemann Hypothesis
is a theorem (Weil, Deligne) and the \'etale fundamental group
provides a computable monodromy group.
Over $\mathbb{Z}$, both ingredients are absent or conjectural.

\subsection{Open questions}
\begin{enumerate}[label=(\arabic*)]
\item \emph{Removing the squarefree hypothesis.}
See Remark~\ref{rem:nonsqfree}.

\item \emph{Uniformity in $d$.}
Does $N_{\mathrm{irr}} \ge 1$ for all $q$ and $d$
with $\mathrm{char} > 2d$?

\item \emph{Shorter intervals.}
Minimize $e$ such that $\{f^2 + s : \deg s \le e\}$
contains an irreducible.

\item \emph{Higher powers.}
For $k \ge 3$, study $[f^k, (f+1)^k]$.
\end{enumerate}

% ═══════════════════════════════════════════════════════════════════════
\begin{thebibliography}{99}

\bibitem{AGLV1993}
  V.\,I.~Arnold, S.\,M.~Gusein-Zade, and A.\,N.~Varchenko,
  \textit{Singularities of Differentiable Maps}, Vol.~II,
  Birkh\"auser, 1988.

\bibitem{bank2015}
  E.~Bank, L.~Bary-Soroker, and L.~Rosenzweig,
  ``Prime polynomials in short intervals and in arithmetic
  progressions,''
  \textit{Duke Math.\ J.}\ \textbf{164} (2015), 277--295.

\bibitem{BHP2001}
  R.\,C.~Baker, G.~Harman, and J.~Pintz,
  ``The difference between consecutive primes, II,''
  \textit{Proc.\ London Math.\ Soc.}\ \textbf{83} (2001), 532--562.

\bibitem{barysoroker2012}
  L.~Bary-Soroker,
  ``Dirichlet's theorem for polynomial rings,''
  \textit{Proc.\ Amer.\ Math.\ Soc.}\ \textbf{140} (2012), 731--740.

\bibitem{ChenConductor}
  R.~Chen,
  ``Conductor incompressibility for Frey curves associated to
  prime gaps: rigidity obstructions to the Wiles paradigm in
  additive prime number theory,''
  Zenodo, 2026.
  \url{https://zenodo.org/records/18682375}

\bibitem{ChenLandau}
  R.~Chen,
  ``On Landau's fourth problem: conductor rigidity and
  Sato--Tate equidistribution for the $n^2+1$ family,''
  Zenodo, 2026.
  \url{https://zenodo.org/records/18683712}

\bibitem{ChenAP}
  R.~Chen,
  ``The 2-2 coincidence: conductor rigidity for primes in
  arithmetic progressions and the Bombieri--Vinogradov barrier,''
  Zenodo, 2026.
  \url{https://zenodo.org/records/18684151}

\bibitem{deligne1977}
  P.~Deligne,
  \textit{Cohomologie \'{e}tale (SGA $4\tfrac{1}{2}$)},
  Lecture Notes in Mathematics, Vol.~569, Springer, 1977.

\bibitem{deligne1980}
  P.~Deligne,
  ``La conjecture de Weil.\ II,''
  \textit{Publ.\ Math.\ IH\'ES} \textbf{52} (1980), 137--252.

\bibitem{dixon1996}
  J.\,D.~Dixon and B.~Mortimer,
  \textit{Permutation Groups},
  Graduate Texts in Mathematics, Vol.~163, Springer, 1996.

\bibitem{FJ2008}
  M.\,D.~Fried and M.~Jarden,
  \textit{Field Arithmetic}, 3rd ed.,
  Springer, 2008.

\bibitem{grothendieck1977}
  A.~Grothendieck \textit{et al.},
  \textit{Cohomologie $\ell$-adique et fonctions $L$} (SGA~5),
  Lecture Notes in Mathematics, Vol.~589, Springer, 1977.

\bibitem{katz1988}
  N.\,M.~Katz,
  \textit{Gauss Sums, Kloosterman Sums, and Monodromy Groups},
  Annals of Math.\ Studies, Vol.~116, Princeton Univ.\ Press, 1988.

\bibitem{katz2001}
  N.\,M.~Katz,
  ``Sums of Betti numbers in arbitrary characteristic,''
  \textit{Finite Fields Appl.}\ \textbf{7} (2001), 29--44.

\bibitem{LW1954}
  S.~Lang and A.~Weil,
  ``Number of points of varieties in finite fields,''
  \textit{Amer.\ J.\ Math.}\ \textbf{76} (1954), 819--827.

\bibitem{mueller1993}
  P.~M\"uller,
  ``Primitive monodromy groups of polynomials,''
  Contemp.\ Math.\ \textbf{186}, AMS, 1995, 385--401.

\bibitem{raynaud1965}
  M.~Raynaud,
  ``Caract\'eristique d'Euler--Poincar\'e d'un faisceau et
  cohomologie des vari\'et\'es ab\'eliennes,''
  S\'em.\ Bourbaki, Exp.~286, 1964/65.

\bibitem{ritt1922}
  J.\,F.~Ritt,
  ``Prime and composite polynomials,''
  \textit{Trans.\ Amer.\ Math.\ Soc.}\ \textbf{23} (1922), 51--66.

\bibitem{rosen2002}
  M.~Rosen,
  \textit{Number Theory in Function Fields},
  Graduate Texts in Mathematics, Vol.~210, Springer, 2002.

\bibitem{serre1979}
  J.-P.~Serre,
  \textit{Local Fields},
  Graduate Texts in Mathematics, Vol.~67, Springer, 1979.

\bibitem{serre1992}
  J.-P.~Serre,
  \textit{Topics in Galois Theory}, 2nd ed.,
  A.\,K.\ Peters, 2008.

\bibitem{SGA1}
  A.~Grothendieck,
  \textit{Rev\^etements \'etales et groupe fondamental} (SGA~1),
  Lecture Notes in Mathematics, Vol.~224, Springer, 1971.

\bibitem{SGA7}
  A.~Grothendieck \textit{et al.},
  \textit{Groupes de Monodromie en G\'eom\'etrie Alg\'ebrique}
  (SGA~7), Lecture Notes in Mathematics, Vols.~288, 340,
  Springer, 1972--73.

\bibitem{turnwald1995}
  G.~Turnwald,
  ``On Schur's conjecture,''
  \textit{J.\ Austral.\ Math.\ Soc.\ Ser.~A}
  \textbf{58} (1995), 312--357.

\bibitem{weil1948}
  A.~Weil,
  \textit{Sur les courbes alg\'ebriques et les vari\'et\'es qui
  s'en d\'eduisent},
  Actualit\'es Sci.\ Ind., Hermann, 1948.

\bibitem{zannier2008}
  U.~Zannier,
  ``On composite lacunary polynomials and the proof of a
  conjecture of Schinzel,''
  \textit{Invent.\ Math.}\ \textbf{174} (2008), 127--138.

\end{thebibliography}

\end{document}
